\subsection{Complete-head circuits and maximal source types}
\label{app:core-head-circuits}

The following family with one owner factor per source admits a closed classification.
It also separates three possible strengthenings of the source-guard
description: maximal coefficients at the original observers, maximal
occurrence types, and maximal extensions to the common carrier $D_I$.

Fix distinct owners $r_1,\ldots,r_n\in I$, where $2\leq n\leq K$,
head variables $H=\{h_1,\ldots,h_n\}$, and a disjoint finite set $A$
of anchor variables. Choose nonempty contexts $L_i\subseteq A$, and put
$V=H\sqcup A$ and $P_i=(H\setminus\{h_i\})\cup L_i$.
For a valuation $w$, write $E_r(w)=\{v:w_v=\{r\}\}$ and
$Z(w)=\bigcup_{a\in A}w_a$. Define
\begin{equation}\label{core:eq:head-function}
 F(w)=Z(w)\cup
       \bigcup_{i:\,P_i\nsubseteq E_{r_i}(w)}w_{h_i}.
\end{equation}
Thus rule $i$ suppresses its head precisely when all its other heads
and its context are $r_i$ singletons. Its own winner is then already
supplied by $L_i$. The canonical coefficient at $h_i$, for observer
$k\ne r_i$, is $[P_i\nsubseteq E_{r_i}]$; for $k=r_i$ it is $1$.
Every anchor coefficient is $1$. In particular these coefficients
factor by owner: $F$ is \emph{source-separable}. A \emph{pure-owner}
coefficient is its restriction to probes whose nonneutral passive
statuses all have one owner.

\begin{theorem}[Complete-head classification]\label{core:thm:head-circuits}
Put
\[
 \operatorname{bad}(w)\ \Longleftrightarrow\
 w_a=\{r_i\}\quad\text{for every }i\text{ and }a\in L_i.
\]
The greatest core operation below $F$ is
\begin{equation}\label{core:eq:head-minorant}
 R(w)=\begin{cases}Z(w),&\operatorname{bad}(w),\\
                    F(w),&\text{otherwise}.
       \end{cases}
\end{equation}
It has the same pure-owner source coefficients and singleton profiles as
$F$. Moreover
\begin{equation}\label{core:eq:head-classification}
 F\text{ is core}\quad\Longleftrightarrow\quad
 \{r_1,\ldots,r_n\}=I\ \text{ or }\
 L_i\cap L_j\ne\varnothing\text{ for some }i\ne j.
\end{equation}
\end{theorem}
\begin{proof}
For $p$, $a\in L_p$, and an ordered list $J$ of at most two distinct
indices different from $p$, start at $h_p$, attach each $h_j$ in order
using $\star_{r_j}$, and finally attach $a$ using $\star_{r_p}$.
Take the derived union of these words and all anchor projections,
and call it $R_0$. Each word has at most four leaves and lies below $F$.
If rule $p$ fires, its final anchor absorbs. If a listed head's rule
$j$ fires, every preceding head is $\{r_j\}$, so its prefix is already
$\{r_j\}$ and absorbs that head at its attachment. Subsequent heads
are also $\{r_j\}$; the final anchor is an unconditional source.
A suppressed head not in the word contributes nothing.

At a bad valuation every word ends in its absorbing anchor, so $R_0=Z$.
Otherwise choose $p$ and $a\in L_p$ with $w_a\ne\{r_p\}$. Take a winner
$k\in F(w)\setminus Z(w)$ supplied by an unsuppressed $h_i$.
If some $b\in L_i$ is not $\{r_i\}$, the word $h_i\star_{r_i}b$
retains $k$. Otherwise every $L_i$ input is $\{r_i\}$, and some
other head is not $\{r_i\}$. Here $p\ne i$. If $w_{h_p}\ne\{r_i\}$,
use $(h_p\star_{r_i}h_i)\star_{r_p}a$. If $w_{h_p}=\{r_i\}$, choose
$j\notin\{i,p\}$ with $w_{h_j}\ne\{r_i\}$ and use
\[
 ((h_p\star_{r_j}h_j)\star_{r_i}h_i)\star_{r_p}a.
\]
Its first prefix is not $\{r_i\}$, and its final anchor is not
$\{r_p\}$, so it retains $k$. Thus $R_0=R$.
Overlapping contexts make bad valuations impossible. If the head owners
exhaust $I$, a bad valuation has $Z=I$. These give the forward
constructions in \eqref{core:eq:head-classification}.

The converse and greatestness use a uniform observer invariant.
Suppose the contexts are pairwise disjoint and $\omega$ is outside the
head owners; it may be an old owner or a newly added observer.
In the ambient winner set $J$ put $N=J\setminus\{\omega\}$.
At row $u_i$ give $h_i$ the menu $J$, every other head and every
anchor in $L_i$ the menu $\{r_i\}$, and all other anchors $N$.
At row $q$ give all heads $J$, each $L_i$ its atom $\{r_i\}$, and
unused anchors $N$. For any word put $U_i=t(u_i)$ and $Q=t(q)$.
Then $U_i\in\{\{r_i\},N,J\}$, and the following hold simultaneously:
\begin{align}
 U_i=\{r_i\},\ U_j=\{r_j\},\ i\ne j
 &\ \Longrightarrow\ \omega\in U_\ell
       \text{ for some }\ell\notin\{i,j\};\tag{A}\label{core:eq:head-invariant-a}\\
 (\forall\ell\ \omega\notin U_\ell),\ U_i=\{r_i\}
 &\ \Longrightarrow\ Q=\{r_i\};\tag{B}\label{core:eq:head-invariant-b}\\
 \omega\in Q&\ \Longrightarrow\ \omega\in U_i
       \text{ for some }i.\tag{C}\label{core:eq:head-invariant-c}
\end{align}
Head, context-anchor, and unused-anchor projections satisfy these
implications directly. For the induction, only a root $r_p$ can erase
$\omega$ at $u_p$; it acts as union at every other $u_i$.
For (A), a root $r_i$ or $r_j$ follows the child providing its own
singleton. At any other root both premises pass to both children.
A witness outside $p$ persists; if both children have only witness $p$,
neither can absorb there. For (B), a root $r_i$ follows its singleton
child, whose query singleton absorbs. At a different root $r_p$,
erasing a child's $p$-witness requires a sibling with $U_p=\{r_p\}$
and $U_i=\{r_i\}$. By (A) that sibling has a witness outside $i,p$,
contrary to a negative row that already passed to both children.
Hence all negative rows pass and (B) applies to both children.
Finally, for (C), follow a child with $\omega$ at the query. If all
parent rows omit $\omega$, that child's only possible witness is $p$.
Its erasing sibling satisfies all negative rows and $U_p=\{r_p\}$;
(B) makes the sibling's query $\{r_p\}$, erasing the query as well.
Owners outside the head list act as union throughout.

The function $F$ omits $\omega$ at every $u_i$ but contains it at $q$,
violating (C). This proves the converse in
\eqref{core:eq:head-classification}. If $t\leq F$ and $w$ is bad,
take $k\notin Z(w)$. Every head owner belongs to $Z(w)$, so $k$ is an
outside observer. Enlarge all heads to $I$ and unused anchors to
$I\setminus\{k\}$, keeping each context at its atom. The negative
rows omit $k$ by the bound on $t$, so (C) makes the enlarged query
omit $k$. Monotonicity gives $k\notin t(w)$. Thus $t\leq R$.

Finally, a head-source monochromatic probe cannot be bad: the nonempty
contexts require at least two different atoms. Anchor probes retain
their source in both functions, even when bad. Their pure-owner
coefficients therefore agree. At bad both functions contain at least
two anchor owners, so their singleton profiles also agree.
\end{proof}

For the remaining results suppose $n=K\geq3$ and the contexts are
pairwise disjoint. Then $F$ is core. Its profiles are
\begin{equation}\label{core:eq:head-profiles}
 \Phi_{r_i}(F)(E)=[V\setminus\{h_i\}\subseteq E].
\end{equation}
We distinguish an occurrence's full owner factors from its contribution
to one observer. At a marked occurrence $o$ of variable $v$ define
\[
 B_{o,r}(E)=\bigwedge_{
   \substack{\text{$r$-owned ancestors of }o\\\text{with sibling }s}}
          \neg\Phi_r(s)(E),\qquad E\subseteq V\setminus\{v\}.
\]
The empty conjunction is $1$. These antitone Boolean functions form
the occurrence's \emph{full type}, ordered pointwise owner by owner;
a type is \emph{actual} for $F$
if it occurs in a word realizing $F$. For observer $k$, the source
coefficient is
\begin{equation}\label{core:eq:head-free-coefficient}
 C_{v,k}((E_r))=
 \bigvee_{o:\,\lambda(o)=v}\ \bigwedge_{r\ne k} B_{o,r}(E_r).
\end{equation}
For disjoint status sets this is the ordinary coefficient. Allowing
the $E_r\subseteq V\setminus\{v\}$ to overlap defines its
\emph{independent-status extension}, which can depend on the realizing
word. Pure-owner probes bound every head type by the vector
\begin{equation}\label{core:eq:head-max-type}
 M_{i,r_i}(E)=[P_i\nsubseteq E],\qquad
 M_{i,s}(E)=1\quad(s\ne r_i).
\end{equation}

\begin{theorem}[Maximal coefficients without maximal types]
\label{core:thm:head-types}
For this core operation $F$, no head has a greatest actual full type.
Nevertheless one realizing word simultaneously attains the greatest
independent-status coefficient at every source and old observer:
anchors and own observers have coefficient $1$, and the coefficient
at $(h_i,k)$ for $k\ne r_i$ is $[P_i\nsubseteq E_{r_i}]$.
\end{theorem}
\begin{proof}
Interpret any realizing word with a fresh owner $\omega$. Singleton
profiles depend only on atomic statuses, so pure-owner coefficients
are unchanged. Its output consequently omits $\omega$ at each $u_i$
in the preceding proof. At $q$ all factors in
\eqref{core:eq:head-max-type} hold, since other heads are ambiguous.
An occurrence of that type would supply $\omega$, contradicting (C).
Thus $M_i$ is not actual.

We exhibit two actual types joining to $M_i$. Fix $j\ne i$ and
$a_j\in L_j$. Let $b_i$ be a fold of $P_i$ at an owner other than
$r_i$, so $\Phi_{r_i}(b_i)(E)=[P_i\subseteq E]$.
Start with marked $h_i\star_{r_i}b_i$. For each
$\ell\notin\{i,j\}$ append at owner $r_\ell$ the sibling
$h_i\star_{r_i}h_j$, and finally append at owner $r_j$ the sibling
\[
 (h_i\star_{r_i}h_j)\star_{r_j}a_j.
\]
Call the result $s_i^j$. It lies below $F$: under rule $i$ the initial
part is $\{r_i\}$ and later stages add only that owner and anchors;
under another rule $\ell$ its attachment absorbs, and later stages
again add only that owner and anchors. Rule $j$ is protected by the
final anchor. If no rule fires, $F$ is the entire input union.
The marked type is $M_i$ except for the extra factor
$B_{r_j}(E)=[a_j\notin E]$. Indeed each earlier foreign sibling's
profile requires $h_i$, absent from a marked-source status set, and
the final sibling's $r_j$ profile reduces to $[a_j\in E]$.

We use two elementary preservation facts. If $s\leq F$, a marked
type of $s$ survives unchanged in the derived union of $s$ with an
$F$ realizer. Whenever its $r$ factor holds, a pure-$r$ probe shows
that neither operand is an $r$ singleton. No subtree of the union
implementation is then an $r$ singleton, so no new $r$ filter is added.
Also, joining two $F$ realizers by any primitive owned choice preserves
their marked types: at a probe where a marked factor holds, the sibling
$F$ cannot be that owner's singleton. Thus finitely many types can
be placed in one $F$ realizer.

Restoration makes the displayed types actual. Two choices of distinct
$j\ne i$ have coordinatewise join $M_i$. Since every actual type is
bounded by $M_i$ and $M_i$ is not actual, there is no greatest one.
For observer $k=r_j\ne r_i$, however, its extra $r_j$ factor is
ignored in \eqref{core:eq:head-free-coefficient}, giving the full
canonical coefficient. At observer $r_i$ the safe direct word
$h_i\star_{r_i}a_i$, $a_i\in L_i$, has no foreign factor. Anchor
projections supply their unconditional coefficients. Restoration and
the second preservation fact put these finitely many types in one
word. The pure-owner upper bounds show that its coefficients are greatest.
\end{proof}

There is also a closed greatest extension on the common carrier $D_I$.
For $v\in D_I^V$ put $E_r(v)=\{x:v_x=a_r\}$, and define $G(v)$ as
follows. If \eqref{core:eq:head-profiles} forces an atom, return that
atom. Otherwise return $p$ precisely when an anchor is $p$, or some
head $h_i$ is $p$ and $P_i\nsubseteq E_{r_i}(v)$; return $n$ in the
remaining cases. This is the product completion prescribed by the
pure-owner coefficients. Every extension of an $F$ realizer is bounded
by $G$, by its occurrence factors.

\begin{corollary}[Greatest common-carrier extension]
\label{core:thm:head-completions}
Among the operations on $D_I$ induced by words realizing $F$, the
greatest is
\[
 H_F(v)=\begin{cases}
          \displaystyle\bigvee_{a\in A}v_a,
             &v_a=a_{r_i}\text{ for all }i,\ a\in L_i,\\
          G(v),&\text{otherwise}.
        \end{cases}
\]
The short-spine word in \cref{core:thm:head-circuits} attains $H_F$,
whereas $G$ is not induced by an $F$ realizer.
\end{corollary}
\begin{proof}
The off-bad short-spine coverage proof works for $p$ just as for an
ordinary winner. Atomic outputs are fixed by the profiles. At bad, all
spines absorb at their final anchor, giving the displayed join.
If that join contains $p$, it already is greatest. Otherwise unused
anchors are at most $n$, and distinct context atoms make the anchor
join $n$. Enlarge all heads to $p$ and unused anchors to $n$.
Use the observer quotient from winner sets over $I\cup\{\omega\}$
to $D_I$: old singletons map to their atoms, other menus omitting
$\omega$ map to $n$, and menus containing $\omega$ map to $p$.
It preserves union and every old-owned choice. Lift the inputs
$a_r,n,p$ to $\{r\},I,I\cup\{\omega\}$, respectively, and project
the resulting word values through this quotient.
At the negative rows pure-owner coefficients forbid $\omega$.
Invariant (C) therefore forbids $p$ at the enlarged query for every
$F$ realizer. Monotonicity gives the same bound before enlargement.
At the query itself $G=p$, proving its nonattainment.
\end{proof}

Already three heads with disjoint singleton contexts give a core
source-separable operation with all these distinctions. Maximal
old-observer coefficients therefore do not imply either a greatest
full occurrence type or attainment of the canonical common-carrier
completion.
